\chapter{Summary and Outlook} 
\label{chap:summary_outlook}

\section{Summary} 
\label{sec:summary_findings} 

This thesis addressed the absence of a systematic approach for verifying the security functions of industrial gateways at the physical IT/OT interface. With Regulation (EU) 2023/1230 making cybersecurity a mandatory part of the conformity assessment of machinery, established functional and safety validation processes lacked a requirement-driven method to verify security functions. To bridge this gap, a prototype methodology was developed using the Design Science Research Methodology. The objective was not to certify a product, but to develop and validate a systematic methodology for verifying security functions. This methodology derives required tests along a traceable chain, mapping the obligations of the regulation through the draft standard prEN 50742 to the component requirements of IEC 62443-4-2. Each requirement is linked to the affected assets, a STRIDE threat, and the treating security function, with test cases specified on both a product-neutral scenario level and a concrete execution level. Furthermore, the approach reuses the existing company test framework through a tool-to-target applicability assessment and classifies each test using the verification taxonomy of IEC 62443-4-1 Practice SVV. 

The central steps of the methodology were successfully demonstrated end-to-end on a real Safe Remote I/O device in an isolated laboratory under the direct-access network position NV1. The derivation yielded fourteen security requirements from Annex III of the regulation and decomposed the device into nine security-relevant assets, establishing a traceable derivation chain that answers the first research question. Answering the second research question, the applicability assessment classified ten company tests before execution, identifying five as applicable, four as non-applicable, and one as blocked. This prevented tools from being run against interfaces the device does not expose. During execution, fifteen test activities produced usable evidence for the HTTP interface, the identification surface, the intervention log, and application-layer availability. Notably, two requirements were assessed as failed due to evidenced device deviations: the intervention log lacked an absolute timestamp and was cleared upon a cold restart, thus failing time-attribution and retention obligations. Additionally, a structured fuzzing run surfaced a temporary, self-recovering disruption of the HTTP service. While the NV1 setup provided substantial evidence (answering the third research question in part), properties depending on cyclic safety communication require passive-observation (NV2) or in-path (NV3) variants, the validation of which remains future work.

Ultimately, the thesis yields significant scientific and practical contributions. Scientifically, it provides a requirement-derivation model translating legal obligations into component-level specifications, a traceability chain linking requirements to concrete evidence, a tool-to-target applicability assessment, and a network-access-variant model clarifying test feasibility. Most importantly, it demonstrates empirically that a component can enter a safe state after a fault while still failing to satisfy cybersecurity requirements, establishing a conceptual distinction between failing safely and being secure. For the industrial partner, the methodology closes the process gap between regulatory expectations and functional testing, yielding a reusable schema that limits additional tooling investments and reduces the effort required for future product releases. Economically, this systematic verification helps avoid the severe costs associated with security-related field incidents or recalls. Ecologically, the early detection of weaknesses extends the safe service life of installed devices, preventing premature hardware replacement. From a social perspective, since the Machinery Regulation treats cybersecurity as a prerequisite for functional safety, this methodology directly contributes to the protection of human operators, fulfilling the regulation's primary social objective.


\section{Future Work} 
\label{sec:future_work}

To advance the current findings, the most immediate step is to realise the two remaining network access variants. Implementing the passive-observation variant NV2 using a mirroring switch or a network tap, and the in-path variant NV3 using a transparent bridge, would allow the execution of scenarios involving safety-telegram integrity, in-path manipulation, watchdog-passivation, and domain-separation. This is crucial for assessing requirements that depend on cyclic safety communication, in particular RQ-002 and RQ-005. Concurrently, outstanding activities requiring no additional variant should be completed within the same campaign, including the induced software-installation event for RQ-008, the configuration change for RQ-009, and the TIA-based identification and maintenance reads for RQ-004, RQ-006, RQ-007, and RQ-013. Together, these steps close the path from the current partial coverage to a complete requirement assessment. Furthermore, applying the methodology to a device that exposes an authenticated and encrypted interface, such as over HTTPS, would additionally enable credential, transport-security, and access-control tests that were not applicable to the current device.

Beyond these immediate testing goals, the methodology itself can be extended. Because prEN 50742 grades the required test depth through the Safety-Related Security Level, the test catalogue can be expanded so that test strictness scales with the level of the affected asset and requirement—ranging, for example, from a simple checksum check at a low level to cryptographic verification under active manipulation at a high level. Another promising direction is the integration of machine learning. A profiling module that learns the normal network behaviour of a device in the testbed could significantly support the detection of anomalous behaviour during a test \cite[pp. 12]{siboni2016securitytestbedinternetthings}.

Finally, to ensure long-term value, the methodology should be permanently embedded in the development and quality-management processes. This includes assigning the responsibilities of the governance model described in Appendix \ref{app:SecureDevelopmentLifecycle}, fixing the point in the lifecycle at which tests are generated and evaluated, and integrating execution into a continuous integration and delivery pipeline for automated security regression testing. It also requires maintaining the traceability matrix when the normative basis changes, such as when prEN 50742 is published in its final form, ensuring the methodology can reliably serve as conformity documentation under the Machinery Regulation. A related open point that requires ongoing attention is the precise regulatory scope for a component that is not itself a machine, particularly since the Safe Remote I/O is a component whose obligations arise primarily through its integration into a larger machine context.

